\begin{Solbook}{1}
 \par \noindent Sắp xếp theo thứ tự tăng dần là: $-22;\ -15;\ 0;\ 2;\ 10$. \hfill \qedEX 
\end{Solbook}
\begin{Solbook}{2}
 \par \noindent  \begin {listEX}[2] \item $268+(-150)+(-168)+250\\=268+(-168)+250+(-150)\\=100+100=200$ \item $5^{1023} \cdot 5^{1000}: 5^{2020}-\left \{2 \cdot \left [152-\left (14-2^{2} \cdot 3\right )\right ]\right \}\\=5^{3}-{2.(152-2)}\\=125-2.150=125-300=-175$ \end {listEX}  \hfill \qedEX 
\end{Solbook}
\begin{Solbook}{3}
 \par \noindent  \begin {listEX}[2] \item $18+5\cdot x =93$\\ \hspace *{1 cm}$\begin {aligned} 5\cdot x&=93-18\\ 5\cdot x&=75 \\ x&=15 \end {aligned}$ \item $x+15 =156-(-30)$\\ $\begin {aligned} x+15&=186\\ x&=186-15 \\ x&=171\end {aligned}$ \end {listEX}  \hfill \qedEX 
\end{Solbook}
\begin{Solbook}{4}
 \par \noindent Gọi hai số cần tìm là $a, b$ ta có tích hai số đó là: $S_1=a \cdot b$.\\ Nếu bớt $16$ đơn vị vào số thứ hai thì tích hai số đó là: $S_2=a \cdot (b-16)=a b-16\cdot a$.\\ Tích mới giảm là: $1476-820=656$\\ Ta có: $S_1-S_2=16\cdot a=656 \Rightarrow a=656: 16=41$\\ $b=1476: 41=36$\\ Vậy số cần tìm thứ nhất là $41$, số thứ hai là $36$. \hfill \qedEX 
\end{Solbook}
\begin{Solbook}{5}
 \par \noindent Gọi $x$ (học sinh) là học sinh khối 6 ($x\in \mathbb {N}$) và $500 < x < 600$)\\ Theo đề ta có: $x \, \vdots \, 18$; $x\, \vdots \, 27$; $x\, \vdots \, 30$\\ Suy ra $x \in BC(18;27;30)$\\ Phân tích mỗi số ra thừa số nguyên tố $18 = 3^2\cdot 2; \quad 27 = 3^3; \quad 30 = 3\cdot 2\cdot 5$\\ BCNN$(18;27;30)=3^3.2.5=270$\\ BC$(18;27;30)=\{0;270;540;810;...\}$\\ Mà $500 < x < 600$ nên $x=540$\\ Vậy khối 6 có $540$ học sinh. \hfill \qedEX 
\end{Solbook}
\begin{Solbook}{6}
 \par \noindent a) Diện tích đất trồng hoa mẫu đơn là $25\cdot 15=375$ (m$^2$) \\ \par b) Số tiền bác An bỏ ra để trồng hoa là \\ $375\cdot 220\,000 + 15\cdot 15\cdot 90\,000+\frac {15\cdot 20}{2}\cdot 160\,000=126\,750\,000$ (đồng) \hfill \qedEX 
\end{Solbook}
\begin{Solbook}{7}
 \par \noindent  a) \begin {center} \begin {tabular}{|l|l|l|l|l|l|} \hline Điểm số & 6 & 7 & 8 & 9 & 10 \\ \hline Số học sinh đạt được & 4 & 4 & 2 & 8 & 12 \\ \hline \end {tabular} \end {center} b) Lớp 6A có 30 học sinh, số học sinh đạt điểm giỏi có $2+8+12=22$ học sinh \hfill \qedEX 
\end{Solbook}
